
\documentclass[UTF8,12pt,a4paper]{ctexart}
\usepackage[margin=2.2cm]{geometry}
\usepackage{amsmath,amssymb,mathtools}
\usepackage{fontspec}
\usepackage{xcolor}
\usepackage{enumitem}
\usepackage{hyperref}
\usepackage{fancyvrb}
\usepackage{fvextra}
\usepackage{microtype}
\usepackage{titlesec}
\usepackage{fancyhdr}
\usepackage{setspace}

\hypersetup{
  colorlinks=true,
  linkcolor=blue!55!black,
  urlcolor=blue!55!black
}

\setstretch{1.18}
\setlist[itemize]{leftmargin=1.8em,itemsep=0.25em,topsep=0.3em}
\setlist[enumerate]{leftmargin=2.1em,itemsep=0.25em,topsep=0.3em}
\titleformat{\section}{\Large\bfseries}{\thesection}{0.8em}{}
\titleformat{\subsection}{\large\bfseries}{\thesubsection}{0.8em}{}
\titleformat{\subsubsection}{\normalsize\bfseries}{\thesubsubsection}{0.8em}{}
\pagestyle{fancy}
\fancyhf{}
\fancyfoot[C]{\thepage}
\renewcommand{\headrulewidth}{0pt}

\DefineVerbatimEnvironment{GuideVerbatim}{Verbatim}{
  breaklines=true,
  breakanywhere=true,
  fontsize=\small,
  frame=single,
  framesep=3mm,
  rulecolor=\color{gray!35},
  bgcolor=gray!7
}

\newcommand{\GuideInlineCode}[1]{\texttt{\detokenize{#1}}}

\begin{document}

\section*{路线 7：Learning-Guided Optimization / 神经网络辅助混合优化}

这条路线回答一个很现实的问题：既然比赛给了 4 张 64GB GPU，我们能不能把量子优化和深度学习结合起来？

可以，但要放在正确位置上。

神经网络不是“直接吐出最优解”的保证器。更靠谱的做法是：让神经网络学习优化流程里的局部策略，然后把它接到 QUBO、退火、QAOA、constrained mixer 和 hybrid workflow 里。

\subsection*{先说结论}

目前我们的体系已经很完整，但不是终点。

原来的六条路线覆盖了：

\begin{itemize}
\item 如何把问题变成 QUBO / Ising。
\item 如何处理约束。
\item 如何用退火采样。
\item 如何用 QAOA。
\item 如何用 constrained mixer 保持部分可行性。
\item 如何用 hybrid 分解处理大规模 MILP / MIQP。
\end{itemize}

新增第七路线后，体系多了一层：

\begin{GuideVerbatim}
学习哪些变量重要、哪些变量该固定、从哪里开始搜索、怎么修复非法解、怎么调量子/退火参数。
\end{GuideVerbatim}

这层非常适合 GPU。

\subsection*{神经网络能学什么}

\subsubsection*{Warm-start}

给一个 QUBO 或 MILP，神经网络预测每个 binary variable 取 1 的概率。

例如：

\begin{GuideVerbatim}
x0: 0.91
x1: 0.08
x2: 0.62
x3: 0.49
\end{GuideVerbatim}

高置信度变量可以作为初始解、固定变量或 QAOA 初态权重。

\subsubsection*{Variable fixing}

如果模型非常确信 \texttt{x0 = 1}、\texttt{x1 = 0}，可以先固定它们，把剩下不确定变量交给 QUBO 子问题。

这和 hybrid 路线天然相连：

\begin{GuideVerbatim}
relaxation / neural policy
  -> fixed variables
  -> smaller QUBO
  -> SA / QAOA / constrained mixer
\end{GuideVerbatim}

\subsubsection*{Branching}

MILP solver 的 branch-and-bound 需要选择“先分裂哪个变量”。GNN 可以学习 branching score。

这不直接求解原问题，但能改变搜索树大小。

\subsubsection*{Repair policy}

当候选解非法时，repair 要决定删哪个、换哪个、翻哪个 bit。

神经网络或强化学习可以学习：

\begin{GuideVerbatim}
当前 violation 状态 -> 下一步 repair action
\end{GuideVerbatim}

\subsubsection*{Parameter scheduling}

QAOA 的 \texttt{gamma/beta}、SA 的温度表、constrained mixer 的初始态权重，都可以由学习模型给出初值或候选分布。

\subsection*{神经网络不能保证什么}

它不能保证全局最优。

深度学习模型通常是启发式。它能利用训练分布里的结构，但如果赛题分布变化很大，它可能失效。

所以比赛方案里必须保留：

\begin{itemize}
\item exact solver 小规模验证。
\item simulated annealing / greedy baseline。
\item feasibility checker。
\item repair。
\item objective recomputation。
\item benchmark gap。
\end{itemize}

神经网络负责“更快找到好解”，不是负责“证明这就是最优解”。

\subsection*{QUBO 怎么变成神经网络输入}

QUBO 本身就是图：

\begin{itemize}
\item 节点：binary variables。
\item 边：二次项 \texttt{q\_ij x\_i x\_j}。
\item 节点特征：线性项、度数、正负耦合和、约束来源。
\item 边特征：coupler coefficient。
\end{itemize}

我们的代码现在提供：

\begin{GuideVerbatim}
QuboGraphFeatureExtractor
\end{GuideVerbatim}

它会导出：

\begin{GuideVerbatim}
node_features
edges
labels
\end{GuideVerbatim}

这就是后续训练 GNN 的数据接口。

\subsection*{训练标签从哪里来}

小规模问题：

\begin{GuideVerbatim}
exact solver -> optimal bitstring
\end{GuideVerbatim}

中规模问题：

\begin{GuideVerbatim}
SA / hybrid / repair -> best feasible incumbent
\end{GuideVerbatim}

大规模问题：

\begin{GuideVerbatim}
classical solver time-limited incumbent
multi-solver ensemble
human-designed heuristic
\end{GuideVerbatim}

标签不一定完美。我们要记录 label source 和 confidence。

\subsection*{4 张 64GB GPU 怎么用}

第一阶段：批量生成训练数据。

跑很多 problem instances，记录 QUBO graph、约束图、solver 输出、repair trace。

第二阶段：训练 GNN warm-start policy。

输入 QUBO 图，输出每个变量取 1 的概率。

第三阶段：训练 repair / local search policy。

把修复过程看成强化学习：

\begin{GuideVerbatim}
state: 当前 bitstring + violation + objective
action: flip / swap / drop / add
reward: violation 降低 + objective 改善
\end{GuideVerbatim}

第四阶段：和量子路线结合。

用学习模型给：

\begin{itemize}
\item QAOA 初始参数。
\item QAOA warm-start state。
\item constrained mixer 的 feasible initial state。
\item 退火初始解。
\item hybrid 子问题变量集合。
\end{itemize}

\subsection*{和量子算法的关系}

这条路线不是量子算法的替代品，而是量子算法的助推器。

对量子退火：

\begin{GuideVerbatim}
学习生成初态 / 子问题 / penalty 配置 / repair 策略
\end{GuideVerbatim}

对 QAOA：

\begin{GuideVerbatim}
学习 gamma/beta 初值、warm-start state、可行初态、参数迁移
\end{GuideVerbatim}

对 constrained mixer：

\begin{GuideVerbatim}
学习 feasible initial state 和 transition graph 偏好
\end{GuideVerbatim}

对 hybrid：

\begin{GuideVerbatim}
学习变量固定、子问题选择和 repair
\end{GuideVerbatim}

\subsection*{更高级量子算法放在哪里}

有更高级的 QAOA 变体：

\begin{itemize}
\item multi-angle QAOA。
\item counterdiabatic QAOA。
\item digitized-counterdiabatic QAOA。
\item gradually changing unitaries。
\item warm-start QAOA。
\item R-QAOA。
\end{itemize}

但它们大多不是全新入口，而是路线 4 和路线 5 的增强。

我们当前最合理的升级路线是：

\begin{GuideVerbatim}
先补 Learning-Guided Optimization
再把 advanced QAOA variants 放进 P2/P3
\end{GuideVerbatim}

因为学习驱动层能同时服务退火、QAOA、constrained mixer 和 hybrid。

\subsection*{代码里的 P0 版本}

当前新增了一个轻量 backend：

\begin{GuideVerbatim}
LearningGuidedSamplerBackend
\end{GuideVerbatim}

它做五件事：

\begin{GuideVerbatim}
1. 从 QUBO 提取图特征
2. 用 linear warm-start policy 给变量打分
3. 先生成 logical candidate，再自动填 slack / auxiliary bits
4. 结合变量固定计划生成候选 bitstring
5. 做局部翻转改进，再交给统一 postprocessor
\end{GuideVerbatim}

这不是最终神经网络，但它把接口打通了。

后续可以把 \texttt{LinearWarmStartPolicy} 替换成 PyTorch/GNN 模型。

现在代码里还有两个接口专门服务后续 GPU 训练：

\begin{GuideVerbatim}
LearningGuidedDatasetBuilder
  -> 把 OptimizationProblem 批量转成 QUBO graph JSONL
  -> 小规模用 exact solver 打标签
  -> 中规模可替换成 SA / hybrid incumbent 打伪标签

VariableFixingPlan
  -> 从 policy 概率和置信度生成 fixed bits / free bits
  -> fixed bits 可交给 hybrid fix-and-optimize
  -> free bits 可继续交给 SA / QAOA / constrained mixer
\end{GuideVerbatim}

也就是说，route 7 的 P0 代码已经不是一句“以后训练 GNN”，而是有真实数据出口和求解入口。

\subsection*{比赛时怎么讲}

可以这样讲：

\begin{GuideVerbatim}
我们的系统不是只跑一个量子算法。
它先把问题转成统一的 QUBO/Ising 表示，再用约束层保证可行性。
求解层有 simulated annealing、QAOA、constrained mixer 和 hybrid decomposition。
在 GPU 可用时，我们再训练 learning-guided policy，为这些求解器提供 warm-start、变量固定和 repair 策略。
\end{GuideVerbatim}

这比“我们用深度学习直接求最优”稳很多，也更像真实工程系统。

\subsection*{共读任务}

读完这份后，回答三个问题：

\begin{enumerate}
\item 为什么神经网络不应该直接承诺最优性？
\item 如果只允许你用 GPU 做一件事，你会训练 warm-start policy、repair policy，还是 QAOA 参数预测器？
\item 对一个大规模赛题，哪些变量适合固定，哪些变量应该留给 QUBO / QAOA 子问题？
\end{enumerate}

能答清这三个问题，你就知道 GPU 和量子优化应该怎么接了。

\subsection*{延伸阅读}

\begin{itemize}
\item \href{https://arxiv.org/abs/1811.06128}{Bengio et al., ML for Combinatorial Optimization}
\item \href{https://arxiv.org/abs/2012.13349}{Nair et al., Neural Networks for MIP}
\item \href{https://arxiv.org/abs/2404.06492}{Darvariu et al., Graph RL for CO}
\item \href{https://arxiv.org/abs/2406.02872}{Liu et al., Automated GNN for CO}
\item \href{https://arxiv.org/abs/2107.02789}{Chandarana et al., Digitized-counterdiabatic QAOA}
\item \href{https://arxiv.org/abs/2605.06858}{Falla and Safro, Constrained Counterdiabatic QAOA}
\end{itemize}

\end{document}
